\documentclass[12pt]{article}
\usepackage[utf8]{inputenc}
\usepackage{url}
\usepackage{hyperref}

\title{Portfolio Project Report}
\author{Meriem EJ-Jabiri - GH1053531 \\ Module: B201 Computer Science Lab \\ University: GISMA University of Applied Sciences}
\date{June 26, 2026}

\begin{document}

\maketitle

\begin{center}
\textbf{ABSTRACT}
\end{center}
This project report describes the creation, organization, development, and online publication of a professional computer science portfolio website. The project was completed as part of a university assessment and was designed to simulate a real application process for a working student position in the technology sector. The main objective was to create a website that presents my personal profile, technical skills, curriculum vitae, and practical programming work in a professional and organized manner. The website is hosted online through GitHub Pages and can be accessed by anyone with an internet connection.

The portfolio serves as a digital space where employers, teachers, and other visitors can learn more about my background, education, and abilities. It demonstrates my understanding of basic web technologies and shows that I can organize files, publish content online, and manage a software project from start to finish. In addition to the website itself, this report explains the structure of the project, the design choices that were made, the tools and technologies used during development, and the lessons learned throughout the process.

\newpage

\section*{INTRODUCTION AND PROJECT BACKGROUND}
Today, technology plays an important role in almost every industry. Because of this, students who want to work in information technology, software development, data analysis, or related fields need ways to demonstrate their skills. A traditional paper resume is still useful, but many employers now expect applicants to have an online presence. A portfolio website allows a student to present their abilities in a clear and professional way.

This project was created to meet those expectations. The task was to build a portfolio website that acts as a central location for important information about me and my work. The website includes details about my education, technical skills, projects, and curriculum vitae. It also provides visitors with a simple way to learn more about my academic journey and future career goals.

The project was divided into two main parts. The first part is the website itself, which is publicly available online. The second part is this written report, which explains the planning, development, and publication process. Together, these elements demonstrate both practical and theoretical knowledge.

One important goal of the project was to understand how websites are deployed online. Publishing a website is different from simply creating files on a computer. It requires hosting services, version control systems, and careful organization. Through this project, I learned how these technologies work together to make a website available to users around the world.

% Added Project Links as strictly required by the assessment brief criteria
\section*{PROJECT LINKS}
As required by the assessment guidelines, the live public links for this project are provided below:
\begin{itemize}
    \item \textbf{Live Portfolio Webpage:} \url{https://Mimz-jb.github.io/Meriem-Ej-Jabiri-portfolio/}
    \item \textbf{GitHub Repository Link:} \url{https://github.com/Mimz-jb/Meriem-Ej-Jabiri-portfolio}
\end{itemize}

\section*{PORTFOLIO STRUCTURE AND ORGANIZATION}
Good organization is essential in every technology project. A well-structured project is easier to maintain, understand, and improve. For this reason, I spent time planning the folder structure before creating the website.

The homepage, stored in the index.html file, acts as the main entry point for visitors. When someone opens the website link, this file is automatically displayed. It contains information about me, my technical skills, and links to additional resources.

A separate folder called cv was created to store my curriculum vitae. Keeping the CV in its own folder helps maintain a clean structure and makes it easier for visitors to download the document. It also prevents important files from being mixed with source code.

Another folder called projects contains practical programming work. This folder currently includes a Python grade calculator program. Separating project files from the homepage keeps the repository organized and allows visitors to quickly locate examples of my programming skills.

This folder structure follows good professional practice. Large software projects often contain many files and directories. Learning to organize information from the beginning is an important skill that will be useful in future academic and professional work.

\section*{DESIGN DECISIONS}
When creating the website, I decided to focus on simplicity and clarity. Many modern websites contain complex animations, large images, and advanced visual effects. While these features can look attractive, they are not always necessary for a student portfolio.

My target audience includes employers, recruiters, lecturers, and assessment graders. These visitors usually want to find information quickly. Therefore, I created a layout that highlights the most important content without unnecessary distractions.

The homepage contains clear sections with headings such as About Me, Technical Skills, Resume, and Projects. These sections make navigation easier and help visitors locate information in a short amount of time.

Another reason for choosing a simple design was performance. Websites that contain large media files often load slowly, especially on older devices or slower internet connections. By using mostly text and lightweight code, the website loads quickly and provides a smooth user experience.

Responsiveness was also considered during development. People may visit the website using desktop computers, laptops, tablets, or smartphones. The layout was designed to adapt to different screen sizes so that the content remains readable regardless of the device being used.

Typography was another important decision. Instead of relying on custom fonts that require additional downloads, I used common system fonts. This improves compatibility and ensures that the website appears consistent across different operating systems.

\section*{TOOLS AND TECHNOLOGIES USED}
Several technologies were used during the creation of this project. Each one played an important role in the development process.

\begin{itemize}
    \item HTML was used to create the structure of the website. HTML defines the different sections, headings, paragraphs, and links that appear on the page. It forms the foundation of every website.
    \item CSS was used to improve the visual presentation. Through CSS, I controlled spacing, margins, text sizes, and overall page appearance. Although the design remains simple, CSS helps make the website more professional and easier to read.
    \item Python was used to create the grade calculation program included in the project folder. Python is one of the most popular programming languages in the world because it is easy to learn and very versatile. The grade calculator demonstrates basic programming concepts such as variables, loops, conditions, and calculations.
    \item IDLE was used as the development environment for writing and testing Python code. It allowed me to run programs, identify errors, and make corrections before uploading files to the repository.
    \item Git was used as a version control system. Version control is extremely important because it records changes made to project files over time. If mistakes occur, previous versions can be reviewed or restored.
    \item GitHub was used to store the project repository online. It acts as a secure cloud platform where project files can be uploaded, managed, and shared.
    \item GitHub Pages was used to publish the website. This service automatically converts repository files into a live website. As a result, visitors can access the project using a web browser without downloading any files.
    \item LaTeX and Overleaf were used to prepare professional documents. These tools help create clean reports and CVs with consistent formatting and professional layouts.
\end{itemize}

\section*{DEVELOPMENT PROCESS}
The development process began with planning. Before writing code, I identified the content that would appear on the website. This included personal information, technical skills, project descriptions, and the CV download section.

After planning the content, I created the basic HTML structure. The initial version focused on functionality rather than appearance. Once the content was organized, CSS was added to improve readability and layout.

The next step involved testing. I opened the website in different browsers to ensure that all sections displayed correctly. Links were checked to verify that visitors could access documents and project files.

At the same time, the Python grade calculator was developed and tested separately. Different test values were entered to confirm that calculations were correct and that the program behaved as expected.

After local testing was complete, Git was used to track changes and create commits. Each major update was saved so that progress could be monitored.

The repository was then uploaded to GitHub. Finally, GitHub Pages was enabled to publish the project online. After deployment, additional testing was performed to confirm that the live website worked correctly.

\section*{BENEFITS OF THE PORTFOLIO WEBSITE}
This portfolio provides several advantages. First, it creates a professional online identity. Employers can learn more about my skills and projects before an interview.

Second, it demonstrates practical abilities. Instead of only describing technical knowledge, the website shows actual examples of work. This makes it easier for visitors to evaluate my capabilities.

Third, the portfolio can continue to grow over time. New projects, certifications, and achievements can be added whenever necessary. As my experience increases, the website can evolve to reflect that progress.

Another advantage is accessibility. Because the website is hosted online, it can be viewed from almost anywhere in the world. This makes it a useful tool when applying for internships, working student positions, or future employment opportunities.

\sectionsection*{PROJECT CHALLENGES}
Although the project was successful, several challenges were encountered during development. One challenge was learning how different technologies interact. HTML, CSS, Git, GitHub, and hosting platforms each have their own rules and workflows.

Another challenge involved troubleshooting. Small mistakes in code can sometimes prevent a website from displaying correctly. Careful testing and debugging were necessary to identify and solve these issues.

Understanding version control was also initially difficult. However, through practice, I learned how commits, repositories, and synchronization work together.

Time management was another important factor. Balancing planning, coding, testing, documentation, and deployment required good organization and discipline.

\sectionsection*{PROJECT REFLECTION AND SELF-EVALUATION}
Reflecting on completed work is an important part of learning. Looking back at this project, I am proud of the final result and the skills developed during the process.

One of the strongest aspects of the project is its organization. The folder structure is clear, logical, and easy to navigate. Visitors can quickly locate the homepage, CV, and project files.

Another strength is simplicity. The website focuses on content rather than unnecessary visual effects. This makes it professional and accessible.

The Python project is also a positive feature because it demonstrates practical programming knowledge. Including a working program helps show that I can apply concepts learned during my studies.

Despite these strengths, there are areas for improvement. The website is currently static, meaning that content must be updated manually. More advanced features could make the site interactive and dynamic.

The visual design could also be improved. While simplicity is valuable, adding modern layouts and navigation elements could make the website look even more professional.

\sectionsection*{FUTURE IMPROVEMENTS}
There are several improvements I would like to make in the future.
\begin{itemize}
    \item The first improvement is adding a database system. A database would allow information to be stored permanently rather than only during program execution.
    \item The second improvement is learning JavaScript. This would allow me to create interactive web applications directly inside the portfolio website.
    \item The third improvement is studying modern web frameworks. Frameworks can help developers build more advanced websites efficiently while maintaining good performance.
    \item Another future goal is to include additional projects. As I complete more coursework and personal projects, I will expand the portfolio to showcase a wider range of skills.
    \item I would also like to improve accessibility by ensuring that the website can be used comfortably by people with different needs and devices.
    \item Finally, I plan to continue learning about cybersecurity, software engineering, and cloud technologies. These areas are becoming increasingly important in the technology industry and would add value to future versions of the portfolio.
\end{itemize}

\section*{ADDITIONAL DISCUSSION AND LEARNING EXPERIENCE}
An additional benefit of creating a portfolio website is the opportunity to practice communication skills alongside technical abilities. Students often focus mainly on programming, but employers also value the ability to explain ideas clearly. Writing descriptions, organizing information, and presenting achievements in a structured format helps develop professional communication. These skills are useful during interviews, team meetings, project presentations, academic reports, internship applications, future employment opportunities, and many other situations.

During the planning phase, I also considered how visitors would interact with the content. Before building any pages, I created a simple outline showing the main sections and their relationships. This approach reduced confusion later because the overall structure had already been defined. Planning is sometimes overlooked by beginners, yet it can save significant amounts of time by preventing unnecessary changes, repeated work, avoidable errors, and organizational problems during development.

Research was another important part of the project. Before making design decisions, I reviewed examples of student portfolios available online. Observing different layouts, navigation methods, and content arrangements helped me understand common practices. Some websites were highly visual, while others focused on written information. Comparing these examples allowed me to identify ideas that matched my goals, audience, and technical experience level at the time of development and initial project creation.

Documentation became an essential activity throughout the project. Keeping notes about decisions, challenges, and completed tasks made it easier to write the final report later. Instead of trying to remember every detail at the end, I recorded important information as development progressed. This habit improved accuracy and demonstrated the value of maintaining clear records during technical projects. Good documentation supports collaboration, learning, maintenance, and future improvements.

Testing involved more than checking whether the website opened correctly. I reviewed links, file locations, text formatting, and overall usability. Different browsers may display content in slightly different ways; therefore, verification was necessary. Testing helped identify small issues before publication and increased confidence in the final result. Even simple projects benefit from careful evaluation because quality assurance contributes directly to a better user experience overall.

The portfolio also encourages continuous learning. Because new technologies appear regularly in the computing field, students must update their knowledge over time. A personal website provides motivation to keep improving skills and adding new accomplishments. Each update represents progress and creates a visible record of development. This can be valuable when reflecting on academic growth, preparing applications, or demonstrating long-term commitment to professional improvement.

Another lesson learned was the importance of attention to detail. Small mistakes such as incorrect file names, broken links, or formatting inconsistencies can affect the professionalism of a website. Correcting these issues required patience and systematic checking. The experience showed that successful projects depend not only on major features but also on many smaller elements working together correctly, efficiently, and consistently across different situations and usage conditions.

Creating the portfolio also improved my confidence with independent problem solving. When difficulties appeared, I searched documentation, reviewed examples, and experimented with different solutions. This process strengthened my ability to learn autonomously, which is an important skill for technology professionals. Many technical challenges do not have immediate answers; therefore, persistence and curiosity are essential qualities. Developing these habits during university studies can provide benefits throughout an entire career.

The use of GitHub introduced me to collaborative development practices even though this project was completed individually. Many companies and open-source communities rely on repositories, version histories, and structured workflows. Understanding these concepts early provides a strong foundation for future teamwork. Experience with version control also reduces the risk of losing work because previous versions remain available when changes need to be reviewed or restored later.

Accessibility remains an important consideration for modern websites. Designers should aim to create content that can be used comfortably by a wide range of people. Clear headings, readable text, logical organization, and meaningful links contribute to accessibility. While this project focuses on fundamental features, awareness of inclusive design principles will help guide future improvements and support the creation of better digital experiences for diverse audiences everywhere.

Security is another topic that became more interesting during the project. Although the website contains mainly public information, learning about online publication raised questions regarding safe development practices. Understanding basic security concepts such as protecting accounts, managing passwords, and maintaining software responsibly is important for anyone working with technology. Future projects will provide opportunities to explore these subjects in greater depth and apply stronger protective measures effectively.

Project management skills were developed alongside technical knowledge. Completing the assignment required setting goals, monitoring progress, and meeting deadlines. Breaking the work into smaller tasks made the process more manageable and reduced stress. This approach is widely used in professional environments because it improves organization and helps teams track responsibilities, resources, and milestones. Efficient planning often contributes significantly to successful project outcomes and sustained productivity.

Feedback can play a valuable role in improving any project. Sharing work with classmates, teachers, or other users may reveal strengths and weaknesses that are not immediately obvious. Constructive comments provide opportunities for refinement and learning. In future versions of the portfolio, I would like to gather feedback from multiple sources and use it to guide design decisions, content updates, and functionality enhancements over time for better results.

Finally, this project represents an important step in my academic journey. As a first-year student, I am still developing many technical abilities, yet completing a full website from planning to publication demonstrates meaningful progress. The experience combined creativity, organization, and problem solving while reinforcing concepts learned during coursework. It also provided motivation to continue exploring computer science and pursuing new challenges with enthusiasm, determination, and confidence. Beyond academic requirements, this experience highlighted the value of consistency and dedication. Regular effort applied over several weeks produced a result that would not have been possible through last-minute work alone. The lesson remains valuable for studies.

\sectionsection*{CONCLUSION}
Building this portfolio website was an important learning experience. The project provided practical exposure to web development, programming, version control, cloud hosting, and professional documentation. From the initial planning stage to the final online publication, every step contributed to a deeper understanding of how technology projects are developed and shared. The completed portfolio demonstrates my current skills while also providing a foundation for future growth.

The project successfully meets its objectives by presenting personal information, technical skills, project work, and a downloadable CV in a professional format. It also proves that I can organize files, manage a repository, publish a website, and document my work effectively.

Most importantly, this experience has increased my confidence as a first-year computer science student. It has shown me that even simple projects can provide valuable practical knowledge when completed carefully and professionally. As I continue my studies, I will build upon the skills gained from this project and use them to create more advanced applications and stronger professional portfolios in the future.

\end{document}